\scanpage{88}

% L1–3
\deleted{\unsure{$h_0\,h_1$}} $h_0\ h_1\ \ldots\ h_n$ sections du fibré projectif $X$. $H_0\ H_1\quad H_n$ \deleted{hyperplans} \deleted{\unsure{$X$}}

% L1–6, right
\margin{$(h_i,H_i)$ centre et cocentre \deleted{de} de pseudo réflexion $\unsure{\tau_i|\mathrm T_i}$ dans $X$\\ $\tau_i\tau_j=\tau_j\tau_i$ si $i,j$ non cons.\\ d'où $\unsure{h_i|H_i}$}

% L3–4
\[
\begin{aligned}
\varnothing=X_{-1}\subset X_0=\{s_0\}&\subset X_1\subset\ldots\subset X_n\subset X=X_{n+1}
 & & (X\text{ de rang }n+1)\\
&&&\tau_i X_j=X_j\ \text{si }i\ne j\quad 0\leq i,j\leq n\\
&\Longrightarrow &&\underline\Phi_i\longrightarrow\Phi_i
\end{aligned}
\]

% L5–7
A) Conditions équivalentes

% L6–7, crossed opening
\deleted{a) $\forall\ 0\leq i\leq n$, $\Phi_i$ ne se réduit à un pt b) $\forall\ 0\leq i\leq n$}

% L7–8
a) $\Phi_0$ engendre lin. $X$ en chaque fibre

% L8
a$'$) l'intersection des $Z\in\Phi_n$ est vide (\deleted{\unsure{ds}} \inserted{en} chaque fibre)

% L9
b) $\forall\ 0\leq i\leq n$, $\Phi_i$ non \unsure{réduite|réduit} à un pt

% L10
c) $\forall\ 0\leq i\leq n$ $\unsure{\tau_i|\Phi_i}(X_i)\ne X_i$\qquad $\exists$

% L11–12
d) si $\underline f_{i-1}<\deleted{\underline f_i\ne\underline f'_i}<\underline f_{i+1}$ dans $\underline\Phi_*$, alors $\underline f_i\ne\underline f'_i$ entre $\underline f_{i-1},\underline f_{i+1}$ tels que $f_i\ne f'_i$.

% L13–14
e) Posant $s_1=\tau_0 s_0$, $s_2=\tau_1s_1=\tau_1\tau_0s_0$, \ldots, $s_{n+1}=\tau_ns_n=\tau_n\tau_{n-1}\ldots\tau_0s_0$ les $(s_i)_{0\leq i\leq n+1}$ engendrent \deleted{affinement} \inserted{lin$^{\mathrm t}$ $X$}

% L15–16
e$'$) Condition duale sur $S_n\inserted{=\text{ par déf }X_n}$, $\deleted{\unsure{\tau_n S_n=}}S_{n-1}=\tau_n S_n$, $S_{n-1}=\tau_{n-1}S_{n-1}=\tau_{n-1}\tau_n S_n\ldots$ $S_{-1}=\tau_0S_0=\tau_0\tau_1\ldots\tau_n S_n$, savoir $\bigwedge_i S_i=\varnothing$.

% L17
\deleted{f) L'une des conditions} f) On a l'ens des \deleted{\unsure{ces}} conditions

% L18–20
\[
\left\{\begin{aligned}
\alpha)&\quad \deleted{\unsure{\text{les}}}(h_i)_{0\leq i\leq n}\ \text{libres \inserted{en tte fibre} -- engendre hyperplan }C\\
\alpha')&\quad(H_i)_{0\leq i\leq n}\ \text{libres \inserted{en ttes fibres} -- }\bigcap_j H_i=\{c\},\ c\ \text{section de }X\\
\beta)&\quad h_i\unsure{\notin|\not\subset}H_{i+1}\ \text{en chaque fibre}\quad 0\leq i\leq n-1
\end{aligned}\right.
\]

% L18–25, left box
\margin{\uline{NB} $\alpha)\,\alpha')$ impliquent $\left\{\begin{smallmatrix}h_i\subset H_j\\h_j\subset H_i\end{smallmatrix}\right.$ si $0\leq i,j\leq n$, $i,j$ non consécutifs (ce qui exprime la cond. de commutation des $\tau_i$) $f)$ implique $X_i=\deleted{\unsure{V|\bigvee}}V(s_0,h_0,\ldots,h_{i-1})=\bigwedge(X_n,H_n,\ldots,H_{i+1})$ $X_i$ et $X_j$ se déterminent mutuellement pour $i<j$ par $X_i=\bigwedge(X_j\cap H_j\cap\ldots\cap H_{i+1})$, $X_j=V(X_i,h_i,\ldots,h_{j-1})$}

% L21–25, right, overwritten
\deleted{\unsure{$\gamma)$}} posons \deleted{\unsure{$D$|$D_0$}} $D_0=H_n\cap\ldots\cap H_1$ (droite), on a déjà a) \deleted{$D_0\cap C=\{d\}$, \unsure{$c$ et $d$|$c\in D_0$}} \inserted{$D_0\cap C=\deleted{\unsure{\{d\}}}\ \text{une section }\{d\}$} \deleted{\illegible{2} \unsure{fixe et $d$} \illegible{2} \unsure{de $\tau_0$})}

% L23–25
$\delta)$ $s_0$ section de \deleted{\unsure{$D_0$}} $D_0^{**}=D_0\smallsetminus\{c\cup d\}=D_0-$\deleted{sections de $X$} pts fixes de $\tau_0\mid D_0$

% L26–28
B) Pour $(\tau_i)$ donné, pour qu'il existe loc. (Zar) un drapeau $(X_i)$ satisfaisant aux conditions équivalentes ci-dessus, il f et s. que $(\tau_i)$ satisfasse les conditions $(\alpha,\alpha',\beta)$ de (A f).

% L29–32
C) Sous les conditions de B (i.e. $\alpha,\alpha',\beta$), l'application $(X_i)\longmapsto s_0$ est une bijection de l'ens \inserted{$\operatorname{Dr}(\tau_*)$} des drapeaux $(X_i)$ adaptés à $(\tau_i)$ et non dégénérés, avec l'ens des sections de $\unsure{\delta_0^{**}|D_0^{**}}$. Plus gén., \deleted{\unsure{$i$ donné}}, \inserted{$\forall\ 0\leq i_0\leq n$},

% L32–36
\deleted{bijection de $\operatorname{Dr}(\tau_*)$ sur l'ens} $(X_{\unsure{j|\nu|*}})\longmapsto X_{i_0}$ est une bijection de $\operatorname{Dr}(\tau_*)$ sur l'ens des sections de $D_i^{**}$, où $D_i^{**}=D_i\smallsetminus\{c_i\cup d_i\}$, $D_i=$ fibré en droites des $(i+1)$-espaces lin. compris entre $Y_{i-1}=V(h_0,\ldots,h_{i-1})$ et $Z_{i+1}=\bigwedge(\deleted{\unsure{X|H_n}}H_n,\deleted{\unsure{\Phi}}\ldots,H_{i+1})$ (resp. de dim $i-1$ et $i+1$), $c_i$, $d_i$ étant les sections de $D_i$ définies resp. par

% L37–39, bottom crossed block
\deleted{$Y''_i=\unsure{Y_{i-1}|Y_{i-1}\vee Z_{i+1}}$, $Z_{i+1}$ $Y'_i$ et $Z'_i=\bigvee(Z_i,Y_{i-1})$ $Y_{i-1}\subset Y'_i,Z'_i\subset Z_{i+1}$ [\uline{NB} $Y_i$ stable par $\tau_0\ldots\tau_i$, $Z_i$ stable par $\tau_i,\ldots,\tau_n$].}

% L38–40, bottom interline
\[
\begin{aligned}
X'_i&=\deleted{\unsure{V}}(Y_{i-1}\vee c) &\text{et}\quad &\deleted{\unsure{V|\bigl(}}Y_{i-1}\vee d\deleted{\unsure{\}|\bigr)}}=\overset{\deleted{\unsure{\text{c'est aussi}}}}{X''_i}\\
&\ \Vert &&\ \Vert\\
&\deleted{\unsure{X}}(Z_{i+1}\wedge \unsure{D|C}) && \deleted{\unsure{X}}(Z_{i+1}\wedge D)
\end{aligned}
\]

% L40, bottom right
\margin{où $D$ est défini dualem$^{\mathrm t}$ à $c$ comme $D=V(h_1\ldots h_{n-1},c)$}

% archive sheet number: 87
